\chapter{Stato dell'arte}
\label{ch:stato_arte}

% esempio da analizzare: https://developer.nvidia.com/blog/building-ai-agents-to-automate-software-test-case-creation/
% https://dl.acm.org/doi/10.1145/3805043
%

% NOTA: il secondo collegamento e' il paper di UnLeakedTestBench \cite{huang2025ult},
% cioe' il dataset usato in questo lavoro: e' la fonte primaria del capitolo.
%
% Il filo del capitolo: si parte da cosa rende buono un test, si mostra come la
% ricerca ha provato a generarli automaticamente, e si arriva alla domanda
% "con quali metriche giudichiamo i test prodotti da un LLM, e quali mancano?".
% Traccia estesa di ogni sezione in appunti/stato_dell_arte.md del repository.


\section{Il testing di unità e il costo della scrittura manuale}


Gli unit test sono un passaggio fondamentale dello sviluppo software, che permette agli sviluppatori di controllare che il codice scritto risulti robusto e funzioni secondo le specifiche richieste. Permette quindi di evidenziare potenziali errori nel software prima che venga effettivamente fatto operare, così come durante l'intero ciclo di vita del codice.

La scrittura manuale dei test può essere vista come un grande collo di bottiglia nel processo di sviluppo. Non solo è causa di larghe perdite di tempo per l'effettiva scrittura di test ma, al contempo, richiede la conoscenza profonda del codice sul quale si sta lavorando \cite{schaefer2024testpilot}.

Scrivere un "test di qualità" significa quindi assicurarsi  che questo sia innzanzitutto corretto ed eseguibile e, allo stesso tempo, deve essere in grado di esercitare il codice in modo esaustivo. Questo è reso complicato dal fatto che nelle moderne architetture di programmi interconnessi è presente spesso una logica intricata, numerose dipendenze e svariati formati di input \cite{huang2025ult}. Nel caso manuale di scrittura, implementare una test suite su un numero molto elevato di funzioni e riuscire a mantenere test di qualità alta è quindi un compito estremamente dispendioso.
% - ruolo degli unit test: incongruenze fra specifica e implementazione, regressioni;
%   correggere presto costa meno che correggere tardi;
% - il collo di bottiglia: scrivere test e' lento e ripetitivo (motivazione condivisa
%   dalle introduzioni di \cite{huang2025ult, wang2025testeval, schaefer2024testpilot});
% - cosa vuol dire "test di qualita'": corretto ed eseguibile, esercita il codice,
%   distingue il codice giusto da quello sbagliato, resta leggibile.
%   Sono le quattro famiglie di metriche della sezione 2.5.


\section{Generazione automatica di test: approcci classici}
\label{sec:classici}
Gli approcci classici per la generazione automatica di testing sono il primo passo per l'attuazione di un nuovo standard non basato sul lavoro diretto dell'uomo in questo campo. Quando si genera un test in primis ci si deve chiedere: con quali input "chiamo la funzione". Questo deve essere fatto in modo tale da esplorarla e coprire più casi possibili. Secondo, conoscendo gli input dati, che cosa ci aspettiamo restituisca?

Ci sono due rami principali:

Nel paradigma Search-Based Software Testing (SBST) la generazione dei test è vista come un problema di ottimizzazione. EvoSuite parte da un insieme casuale di test. Viene usata la copertura come fattore di merito: le suite che hanno la "coverage" più alta sopravvivono e vengono combinate tramite operatori genetici come crossover e mutazione, mentre quelle che non soddisfano la specifica desiderata vengono progressivamente scartate \cite{fraser2011evosuite}. Pynguin sfrutta lo stesso paradigma specificamente su Python (laddove EvoSuite è eseguito per classi nel linguaggio Java) nel quale lo strumento, privo di tipizzazione statica, deve effettuare un'analisi preliminare. Viene quindi fatta una scansione sul codice per trovare tipi espliciti oppure si controllano le type annotations e qualora manchino comunque le informazioni necessarie Pynguin sceglierà i tipi in modo casuale \cite{lukasczyk2022pynguin}.

L'esecuzione simbolica, d'altro canto, affronta la generazione dei test come un problema matematico per raggiungere specifici cammini di esecuzione della porzione di codice sulla quale si sta lavorando \cite{king1976symbolic}. Il programma non viene eseguito con valori numerici specifici ma con simboli algebrici che rappresentano valori arbitrari. Durante l'esecuzione (in particolare quando si incontra una diramazione condizionale) viene aggiornata la "Path Condition" che descrive un determinato cammino e gli input necessario a percorrerlo. Risolvendo i vincoli logici che si incontrano è possibile risalire al valore numerico necessario a rispettare la condizione.

Entrambe queste famiglie di strumenti hanno un limite di fondo: seppur eccellendo nell'esplorazione e nel raggiungere determinati rami specifici, non dispongono din un "oracolo": il valore atteso da inserire nell'asserzione viene ottenuto limitandosi a "fotografare" la situazione già esistente senza controllare le specifiche desiderate. Se la funzione contenesse un difetto, il test generato in modo automatico lo farebbe passare come corretto: saper raggiungere il codice e conoscere cosa quel codice debba restituire sono due qualità distinte e la modalità di generazione automatica dei test ha solo la prima.


% - search-based: EvoSuite per Java \cite{fraser2011evosuite}, Pynguin per Python
%   \cite{lukasczyk2022pynguin}: cercano input che massimizzano la copertura;
% - simbolici: esecuzione simbolica \cite{king1976symbolic} e solver SMT
%   \cite{demoura2008z3} per derivare input che raggiungono cammini specifici;
% - il limite comune: test poco leggibili e soprattutto assenza di un oracolo.
%   Sanno raggiungere una riga, ma il valore atteso lo ricavano eseguendo il
%   codice cosi' com'e': fotografano il comportamento attuale invece di
%   verificare quello desiderato.
%
% Questa sezione introduce la distinzione fra "raggiungere il codice" e "sapere
% cosa deve restituire", che ritorna in tutto il resto della tesi.


\section{Generazione di test con modelli linguistici}
Un primo passo nella direzione di unire in un unico workflow sia l'esplorazione del codice che la deduzione di un oracolo è avvenuto grazie allo sviluppo di reti di apprendimento automatico. Un esempio di sviluppo di questo tipo di tecnologia è rappresentato da TOGA (Test Oracle Generation) \cite{dinella2022toga}. TOGA si colloca come ponte tra la generazione automatica e i modelli neurali: delega a due sottosistemi due compiti diversi. Ad uno strumento (come può essere EvoSuite) affida la generazione della sequenza di esecuzione (test prefix), mentre affida ad una rete neurale il compito di dedurre l'oracolo. Questo però presuppone la disposizione di un ampio corpus di addestramento della suddetta rete, ed è questo ciò che i modelli linguistici eliminano.

Il salto nella generazione dei test è avvenuto proprio con i modelli linguistici di grandi dimensioni. Il vantaggio offerto dall'utilizzo di questi strumenti sta nella capacità di avvicinarsi di più al risultato corretto, "indovinando" a partire dal nome della funzione, dalla docstring oppure in altri casi dagli stessi schemi visti in fase di addestramento \cite{schaefer2024testpilot}. Un ulteriore vantaggio è la leggibilità: i nomi dei test sono descrittivi e le asserzioni esprimono il comportamento della funzione. Il valore atteso proviene da un'ipotesi su ciò che il codice dovrebbe fare, non da ciò che fa.
% - il salto: gli LLM producono test simili a quelli scritti dagli sviluppatori,
%   con nomi parlanti e assert espressi in termini di comportamento atteso;
% - primi lavori sugli oracoli: ATLAS \cite{watson2020atlas}, TOGA \cite{dinella2022toga};
% - TestPilot \cite{schaefer2024testpilot}: genera test senza addestramento ne'
%   few-shot, costruendo il prompt con firma, documentazione ed esempi d'uso;
%   70,2% di copertura di istruzioni mediana contro il 51,3% dello stato dell'arte
%   precedente; introduce l'analisi delle assert non banali;
% - approcci iterativi e ibridi: CoverUp \cite{pizzorno2024coverup} guidato dalle
%   righe non coperte, CodaMOSA \cite{lemieux2023codamosa} che usa l'LLM per
%   sbloccare la ricerca quando si arena;
% - la formulazione del prompt incide sulla qualita' dell'output: si veda
%   \cite{tony2025prompt} per l'ottimizzazione automatica dei prompt nella
%   generazione di codice Python;
% - linguaggi a tipizzazione dinamica: in Python mancano i tipi statici che
%   guidano gli strumenti tradizionali.


\section{Benchmark per la valutazione}
\label{sec:bench}
I test generati dagli LLM devono poter essere valutati tramite dei parametri per evidenziare la loro efficacia. Diversi studi si sono prodigati per analizzare come diversi modelli sviluppassero test e su questi fosse possibile estrarre dei risultati su più metriche contemporaneamente.

In UnLeakedTestbench \cite{huang2025ult} sono raccolte funzioni che rispettano a priori diversi attributi che permettono di valutare le capacità del modello sulla singola FUT (Function Under Test). 

Granularità: valutazione a livello di funzione, al contrario di classi o file. Realismo: generare test partendo per esempio da codice attinto da LeetCode equivale ad esercitare il modello su qualcosa di artificiale e non attinente a come verrebbe scritto da un umano. ULT attinge a The Stack v2 \cite{lozhkov2024stackv2}, un'ampia raccolta di codice open-source dalla quale sono state estratte funzioni scritte in linguaggio Python di software reale. Le funzioni candidate inoltre hanno una complessità ciclomatica \cite{mccabe1976complexity} non inferiore a dieci in modo tale da eliminare a priori le funzioni triviali che avrebbero considerevolmente favorito la performance dei modelli sulla suite.
Infine le funzioni prese sotto esame sono esenti da contaminazione: questo significa che sono codice non presente in fase di addestramento del modello in considerazione. Questo infatti avrebbe favorito la "memoria".

Ad ULT è stato affiancato PLT: un superset che in maniera preliminare opera come ULT per gli altri attributi delle funzioni contentute, infatti include lo stesso ULT. All'interno però si trovano anche  tutte quelle funzioni che sono state scartate per contaminazione e risulta utile per accreditare ancor di più la tesi per la quale il codice contaminato falsa i risultati ottenuti.
% Confronto fra i banchi di prova esistenti. Punti da sviluppare:
% - granularita': molti benchmark valutano a livello di file o classe; ULT porta
%   la valutazione a livello di funzione, l'unita' naturale dello unit testing;
% - realismo: il codice LeetCode e' artificiale e gonfia i risultati; ULT attinge
%   a The Stack v2 \cite{lozhkov2024stackv2} e tiene solo funzioni con complessita'
%   ciclomatica \cite{mccabe1976complexity} maggiore o uguale a 10;
% - contaminazione dei dati: se funzioni e test erano nell'addestramento, il
%   modello ricorda invece di generalizzare. ULT filtra le funzioni con test
%   pubblicamente disponibili; PreLeakedTestBench e' il sovrainsieme che le
%   include e serve a misurare l'effetto della contaminazione per differenza;
% - divario fra benchmark: dal 91,79% di Pass@1 su TestEval al 41,32% su ULT.

\begin{table}[H]
\centering
\caption{Confronto fra i principali benchmark per la generazione di unit test.}
\label{tab:benchmark}
\begin{tabular}{lllll}
\hline
\textbf{Benchmark} & \textbf{Granularità} & \textbf{Dimensione} & \textbf{Metriche} & \textbf{Decont.} \\
\hline
TestEval \cite{wang2025testeval} & programma & 210 programmi & copertura & no \\
TestGenEval \cite{jain2025testgeneval} & file & 1.210 coppie & pass@k, cov., mut. & no \\
ULT \cite{huang2025ult} & funzione & 3.909 funzioni & Pass@k, LCov, BCov, Mut & sì \\
\hline
\end{tabular}
\end{table}

% Qui va dichiarato che questo lavoro usa ULT, e perche': livello di funzione,
% codice reale, nessun rischio di contaminazione.


\section{Metriche per la qualità dei test}
La valutazione di un test generato da un modello linguistico va oltre il fatto che questo sia sintatticamente valido e sia eseguito senza errori. Sebbene siano condizioni necessarie, non sono sufficienti a descriverne tutti gli aspetti. La letteratura ne affronta molti altri: quanto il test sia capace di esplorare il codice, se sia in grado di individuare un'implementazione difettosa, la ridondanza e la duplicazione.

Solo la lettura congiunta di queste misure è capace di restituire una panoramica adeguata sull'affidabilità del test.


% Sezione centrale: giustifica il titolo della tesi. Per ognuna: cosa misura,
% come si calcola, cosa non vede. L'ordine e' argomentativo: si parte da cosa
% significa che un test "funziona", si passa a quanto esercita il codice, si
% mostra perche' quella misura da sola inganna, e si arriva alle metriche che
% coprono il vuoto.

\subsection{Correttezza ed eseguibilità}
In primis il test deve essere generato: un modello può fornire una risposta vuota o non rispondere affatto anche se il prompt fornito è lo stesso di quelli per cui ha risposto correttamente. Questo tipo di fallimento avviene quindi nella pipeline di generazione del test e va contato separatamente dalle altre metriche: queste infatti possono essere estratte solo alla fine del processo.
In ordine segue poi la correttezza sintattica del test: questa può essere verificata con un parser AST ancor prima che il test venga eseguito. Poi il codice sintatticamente corretto ma che non arriva all'esecuzione per esempio a causa di import mancanti o incompleti. Infine la differenza tra test eseguiti che falliscono e test che invece passano.

La metrica che esprime la proporzione di test passati è chiamata Pass@k \cite{chen2021humaneval}: test corretti sintatticamente che arrivano a termine e le cui assert esprimono aspettative esatte sul comportamento della funzione. Questa qualità è valutata sul singolo test: è il numero di test validi per funzione diviso il numero di test richiesti ($K$). Il calcolo di Pass@k passa attraverso $C_i$, definito come il numero di test corretti generati per la i-esima funzione (dove $0 \leq C_i \leq K$). Se chiamiamo con $N$ il numero totale di funzioni in un benchmark, secondo la formulazione di Huang et al. \cite{huang2025ult} otteniamo: \begin{equation}
Pass@k = \frac{\sum_{i=1}^{N} C_i}{N \times K}.
\label{eq:passk}
\end{equation} 

Questa metrica nasce per la valutazione dei modelli nella generazione di codice \cite{chen2021humaneval} e viene adottata da tutti i benchmark di testing \cite{huang2025ult, wang2025testeval, jain2025testgeneval}. Non è però capace di distinguere autonomamente un errore di forma da un errore di ragionamento ed è quindi necessario attribuire dei livelli di fallimento.
% - la prima domanda non e' se il test sia buono ma se esista: un modello puo'
%   non rispondere affatto, o restituire una risposta vuota. E' il fallimento
%   della pipeline, e va contato separatamente da tutto il resto;
% - i livelli successivi, in ordine di severita' crescente: codice
%   sintatticamente non valido (verificabile con un parser AST, senza eseguire
%   nulla), codice che non arriva all'esecuzione (import mancanti, errori di
%   raccolta), test eseguiti che falliscono, test che passano;
% - Pass@k \cite{chen2021humaneval}: proporzione di test corretti, cioe' che
%   compilano, arrivano a termine e le cui assert esprimono aspettative valide.
%   Da precisare che la metrica e' definita sul SINGOLO test, non sul file: in
%   \cite{huang2025ult} e' il numero di test corretti per funzione diviso il
%   numero di test richiesti. Un file con sei test di cui quattro passano non e'
%   ne' un successo ne' un fallimento;
% - la metrica nasce nella valutazione dei modelli per la generazione di codice
%   \cite{chen2021humaneval, austin2021mbpp} ed e' adottata da tutti i benchmark
%   di test \cite{huang2025ult, wang2025testeval, jain2025testgeneval};
% - perche' la distinzione conta: un file non eseguibile e un test che fallisce
%   dicono cose diverse sul modello, il primo e' un errore di forma, il secondo
%   di ragionamento.

\subsection{Copertura del codice}
La copertura è un'altra metrica nella valutazione dei test. In linea generale può essere definita come quella qualità che esprime quanto il codice viene esplorato. Secondo la definizione adottata da Huang et al. \cite{huang2025ult} può essere espressa secondo due sottometriche: coprtura di riga (LCov@k) e copertura di ramo (BCov@k). La prima è la percentuale delle linee di codice eseguite dai test case generati dalle prime k richieste (in questo caso k indica il numero di richieste fatte al modello). La seconda è la percentuale di rami di esecuzione attraversati dai casi di test nelle prime k richieste. Data una suite di funzioni ad alta complessità ciclomatica, la Branch Coverage è quella che discrimina meglio tra modelli ed è la più sensibile alla contaminazione dei dati \cite{huang2025ult}.

Una nota fondamentale da riportare è che la coverage è ottenuta anche dai test che falliscono perché un assert errato esegue comunque la funzione prima di sollevare l'eccezione.

% - copertura di istruzioni e di ramo; copertura di cammino nei casi mirati;
%   guadagno incrementale al crescere del numero di test;
% - la copertura di ramo e' piu' severa, e su funzioni ad alta complessita'
%   ciclomatica e' quella che discrimina davvero fra i modelli; secondo
%   \cite{huang2025ult} e' anche la piu' sensibile alla contaminazione dei dati;
% - attenzione al denominatore: media sui soli file eseguibili oppure su tutti i
%   campioni? Sono due domande diverse e vanno dichiarate entrambe;
% - osservazione che prepara la sottosezione seguente: la copertura si ottiene
%   anche dai test che falliscono, perche' un assert sbagliato esegue comunque
%   la funzione prima di sollevare l'errore.

\subsection{I limiti della copertura come indicatore}
L'efficacia della misura di copertura del codice è dibattuta in letteratura. Secondo lo studio effettuato da Inozemtseva e Holmes \cite{inozemtseva2014coverage} la correlazione tra il livello di copertura e l'efficacia nell'individuare i difetti è è più debole di quanto comunemente si assuma. Gli autori hanno descritto come la dimensione della suite di test possa agire come fattore confondente. Una volta reso costante il numero di test presenti nella suite, la correlazione sopra citata scende a bassa o moderata. La critica appare evidente se accostata a ciò che è stato detto nella Sezione~\ref{sec:classici}: raggiungere una riga non implica sapere cosa questa debba produrre.
% - Inozemtseva e Holmes \cite{inozemtseva2014coverage}: controllando il numero
%   di test nella suite, la correlazione fra copertura ed efficacia scende a
%   bassa o moderata; la dimensione della suite e' un fattore confondente;
% - il punto e' quello gia' incontrato nel 2.2 con gli strumenti search-based:
%   raggiungere una riga non significa sapere cosa quella riga debba produrre.
%   La copertura misura il primo aspetto e ignora il secondo;
% - conseguenza per la tesi: la copertura da sola e' un indicatore parziale, e
%   serve affiancarle metriche che guardino l'oracolo. E' l'argomento che motiva
%   l'intero lavoro;
% - (opzionale) studio di replicabilita' recente sulle suite generate da LLM:
%   arXiv:2607.22880.

\subsection{Capacità di rilevare difetti}
Il mutation score (Mut@k) risponde direttamente al limite sopra descritto della copertura nella valutazione. Viene effettuato introducendo piccole modifiche sintattiche (un "+" che diventa "-", un "$>$" che diventa "$\geq$") ottenendo versioni difettose della funzione, dette mutanti. Una test suite uccide un mutante se fallisce quando viene eseguita sul codice mutato, mentre passa correttamente su quello originale. Nel lavoro sviluppato da Huang et al. è stato posto un timeout pari a 120 secondi in modo da evitare che la suite di test elabori per troppo tempo su un mutante \cite{huang2025ult}. Gli stessi autori riportano la formulazione per ottenere il punteggio: \begin{equation}
Mut@k = \frac{\sum_{i=1}^{N} M_i}{N}.
\label{eq:mutk}
\end{equation}, 
dove $N$ è il numero totale di funzioni in un benchmark mentre $M_i$ la proporzione di mutanti uccisi dalla suite generata per la i-esima funzione.

Il mutation score è la metrica più difficile da ingannare: il test per uccidere il mutante deve essere in grado di discernere il codice sano da quello difettoso, quindi deve essere in possesso di un oracolo funzionante. Il costo tuttavia è elevato perché richiede l'intera riesecuzione della test suite per ogni mutante: è quindi spesso riportato a livello di sottoinsiemi.
% E' la risposta diretta al limite appena descritto.
% - mutation score \cite{papadakis2019mutation}: si iniettano difetti sintetici
%   nel codice sotto test (un "+" che diventa "-", un ">" che diventa ">=") e si
%   misura quanti di questi mutanti vengono "uccisi", cioe' fanno fallire la
%   suite. Lo strumento usato da \cite{huang2025ult} per Python e' Cosmic Ray;
% - la validita' dei mutanti come sostituti dei difetti reali e' argomentata in
%   \cite{just2014mutants};
% - e' la metrica piu' difficile da ingannare: per uccidere un mutante il test
%   deve distinguere il codice sano da quello guasto, quindi deve possedere un
%   oracolo corretto e non limitarsi a passare per la riga giusta;
% - costo: richiede di rieseguire l'intera suite su ogni mutante, ed e' la
%   ragione per cui viene spesso riportata su sottoinsiemi.

\subsection{Duplicazione e ridondanza}
La misura della copertura può essere gonfiata dal modello se produce diverse varianti per lo stesso caso: questa particolare situazione ha come effetto che i numeri salgono ma il valore diagnostico no.
A livello operativo è quindi utile riportare in valore percentuale, così da prescindere dalla lunghezza complessiva, le righe che si ripetono fra i test della stessa suite rapportate al totale delle righe di test.

La duplicazione risulta essere una metrica complementare e non sostitutiva della copertura. Questa misura non viene riportata da nessuno dei benchmark della Sezione~\ref{sec:bench}. Il lavoro che gli si avvicina di più è TestPilot \cite{schaefer2024testpilot}, che misura la distanza tra i test generati e quelli già presenti. La differenza è che in questo caso la ridondanza è misurata confrontando con i test già presenti all'esterno (per esempio quelli generati manualmente) ma non internamente alla suite.
% - un modello puo' gonfiare la copertura producendo molte varianti dello stesso
%   caso: i numeri salgono, il valore diagnostico no;
% - definizione operativa: percentuale di righe duplicate sul totale delle righe
%   di test. Riportarla in percentuale la rende indipendente dalla lunghezza dei
%   file e quindi confrontabile fra modelli;
% - misura complementare, non alternativa: due suite con la stessa copertura ma
%   diversa duplicazione non hanno lo stesso valore;
% - nessuno dei benchmark del 2.4 la riporta.

\subsection{Qualità del codice di test}
I test smell sono difetti di progettazione che non fanno fallire i test ma li rendono fragili, difficili da mantenere e riutilizzare. Sono catalogati da van Deursen et al. \cite{vandeursen2001testsmells}. L'idea dietro il lavoro è che il fatto di passare non è garanzia di qualità.

Un difetto ricorrente nei test generati dai LLM è l'Assertion Roulette,
che si verifica quando un test contiene molte asserzioni prive di messaggi esplicativi. Questo è qualcosa di altamente problematico dal punto di vista ingegneristico perché non fornisce informazioni sulla causa di un errore, solo sul fatto che è avvenuto.

Alla stessa area appartengono la leggibilità e la manutenibilità del codice di test: i test generati da un modello linguistico tendono ad essere più leggibili di quelli prodotti dagli strumenti search-based, rispettano una struttura convenzionale e adottano nomi descrittivi.

Nello stesso contesto si colloca l'analisi delle asserzioni banali (trivial assertions), affrontata da Schäfer et al. in TestPilot \cite{schaefer2024testpilot}. Si tratta di asserzioni che non verificano la reale correttezza del comportamento del software, poiché risultano prive di significato o tautologiche pur contribuendo nominalmente alla copertura del codice.

Tutte le dimensioni sopra citate non vengono misurate in questo lavoro perché richiedono valutazione umana o strumenti dedicati, e sono riprese fra gli sviluppi futuri.
% - test smell: difetti ricorrenti di progettazione dei test, catalogati in
%   \cite{vandeursen2001testsmells}. Sui test generati da LLM prevalgono
%   Assertion Roulette (molte assert senza messaggi, non si capisce quale abbia
%   fallito) e Magic Number Test (opzionale: arXiv:2410.10628);
% - leggibilita' e manutenibilita': i test da LLM sono piu' leggibili di quelli
%   search-based, con nomi parlanti e struttura convenzionale, ma restano meno
%   curati di quelli scritti a mano;
% - assert non banali \cite{schaefer2024testpilot}: assert che verificano davvero
%   il comportamento, contro assert vuote o tautologiche;
% - queste dimensioni non vengono misurate in questo lavoro perche' richiedono
%   valutazione umana o strumenti dedicati: vanno dichiarate qui e riprese negli
%   sviluppi futuri.

\subsection{Metriche sul comportamento del modello}
\label{sec:comportamento}
Queste metriche non misurano il test ma il processo che lo produce. Sono tre:
\begin{description}
  \item[Aderenza al formato] ovvero il rispetto dei vincoli espressi nel prompt: import, assenza di markdown, nessuna riscrittura della funzione sotto test;
  \item[Troncamento della risposta per esaurimento del budget di token] un file troncato può superare il controllo sintattico se il taglio cade a fine funzione: sembra sano ma è incompleto. Ecco perché va trattato come categoria a sé e non lasciato dentro "non eseguibile".
  \item[Tempo per generazione] valutabile solo a parità di infrastruttura.
\end{description}
Questa famiglia è assente dai lavori della Sezione~\ref{sec:bench}, che valutano il mero output (il test).
L'aderenza al formato e il troncamento permettono di risalire all'origine del fallimento: c'è distinzione tra non rispettare il prompt fornito e sbagliare il ragionamento.
% Non misurano il test ma il processo che lo produce, e sui modelli piccoli sono
% spesso piu' informative delle metriche classiche: spiegano PERCHE' un campione
% e' finito in un certo ramo, non solo che ci e' finito.
% - aderenza al formato richiesto: rispetto dei vincoli espressi nel prompt
%   (import corretto, numero di test entro il range, assenza di markdown,
%   nessuna riscrittura della funzione sotto test);
% - troncamento della risposta per esaurimento del budget di token, segnalato
%   dal campo finish_reason. Da trattare come categoria a se': un file troncato
%   puo' superare il controllo sintattico se il taglio cade a fine funzione;
% - tempo per generazione, confrontabile solo a parita' di infrastruttura;
% - questa famiglia e' assente dai benchmark del 2.4, che valutano l'output e
%   non il comportamento.

\section{Sintesi e posizionamento del lavoro}
\label{sec:posizionamento}
Dai concetti sopra espressi emergono due elementi principali che lasciano spazio ad un contributo:
\begin{enumerate}
  \item Il primo riguarda i modelli valutati nei benchmark. Nella Sezione~\ref{sec:bench} la maggior parte delle valutazioni di letteratura si focalizza su Modelli Linguistici di Grande Taglia (LLM) o su modelli ampiamente specializzati sul codice. In ULT \cite{huang2025ult} vengono usati per esempio DeepSeekCoder, Qwen2.5-Coder e Phi-4. Restano fuori dall'esplorazione modelli di piccola taglia, comodi per poter essere fatti lavorare in locale su macchine dalle risorse computazionali limitate senza la necessità di ricorrere ad API remote.
  \item Il secondo riguarda l'incompletezza delle metriche sulle quali sono stati valutati i suddetti modelli. Nessuno dei benchmark considerati si occupa di valutare secondo delle metriche aggiuntive come la ridondanza fra i test e quelle sul comportamento del modello spiegate nella Sezione~\ref{sec:comportamento}. Su modelli di dimensioni ridotte come quelli utilizzati i fallimenti legati al mancato rispetto del formato, al troncamento della risposta o alla ripetizione dello stesso caso avvengono con frequenza più elevata.
\end{enumerate}
Questo lavoro si inserisce proprio in questo contesto: valutare l'efficacia contemporaneamente alle limitazioni dei modelli di piccola taglia nell'ambito della generazione dei test.



% Lacune individuate:
% 1. i modelli valutati sono quasi sempre grandi o specializzati sul codice; ULT
%    \cite{huang2025ult} ne valuta 12, ma i generalisti molto piccoli (1--8
%    miliardi di parametri) sono poco esplorati, benche' siano gli unici
%    eseguibili su una macchina ordinaria;
% 2. le metriche vengono riportate in parallelo ma raramente messe in tensione
%    fra loro: il divario fra copertura alta e correttezza bassa merita di essere
%    misurato esplicitamente;
% 3. restano fuori dai benchmark la ridondanza fra i test generati, il numero di
%    riesecuzioni della stessa riga e l'aderenza al formato richiesto.
%
% Contributo della tesi: valutare tre modelli di taglia crescente a parita' di
% prompt, dataset e parametri, aggiungendo alle metriche classiche misure di
% ridondanza e di comportamento, e discutendo cosa ciascuna metrica vede e cosa
% nasconde.
